\pdfvariable trailerid {[<C3332026090300000000000000000000><C3332026090300000000000000000000>]}
\documentclass[10pt]{article}
\usepackage[margin=0.72in]{geometry}
\usepackage{amsmath,amssymb,amsthm,mathrsfs,microtype,hyperref}
\hypersetup{colorlinks=true,linkcolor=blue,urlcolor=blue}
\providecommand{\CRevisionRound}{2}
\newtheorem{theorem}{Theorem}
\newtheorem{lemma}[theorem]{Lemma}
\newtheorem{remark}[theorem]{Remark}
\newcommand{\EE}{\mathbb E}
\newcommand{\one}{\mathbf 1}
\newcommand{\cT}{\mathcal T}
\newcommand{\cS}{\mathcal S}
\setlength{\emergencystretch}{3em}
\title{Relaxed Randomized Gossip on the Complete Graph:\\
Exact Covariance Spectrum and Sharp Consensus Law}
\author{Route-A source-local certificate HCS-C333}
\date{3 September 2026}

\begin{document}
\maketitle
\begin{abstract}
For uniform pairwise relaxed averaging on the complete graph, we derive the
exact first moment and sharp mean-square disagreement law at every finite
time.
\ifnum\CRevisionRound>0
The entire second-moment transfer on $\operatorname{Sym}^2(\one^\perp)$
splits into three explicit orthogonal blocks; all eigenvalues,
multiplicities, projectors, and matrix-valued iterates are closed.
\fi
\ifnum\CRevisionRound>1
This yields the exact statistical covariance, a finite-time tail bound,
almost-sure consensus for $0<\eta<1$, and a complete treatment of
$N=1,2,3$ and $\eta=0,1$.  In particular, $\eta=1$ is a random-transposition
boundary and not a consensus regime.
\fi
\end{abstract}

\section{Frozen stochastic owner}
Fix $N\geq2$ and $0\leq\eta\leq1$.  Independently at every integer time,
choose an unordered pair $\{i,j\}$ uniformly from the edges of $K_N$ and set
\begin{equation}\label{eq:update}
 x_{t+1}=W_{ij}x_t,
 \qquad W_{ij}=I-\eta d_{ij}d_{ij}^{\mathsf T},
 \qquad d_{ij}=e_i-e_j.
\end{equation}
Thus the selected entries become
$(1-\eta)x_i+\eta x_j$ and $\eta x_i+(1-\eta)x_j$.
The endpoint $\eta=1/2$ is ordinary pairwise averaging, while $\eta=1$
merely swaps the two entries.  The initial vector is deterministic.

Put
\[
 P=I-\frac1N\one\one^{\mathsf T},\qquad
 \bar x=\frac1N\one^{\mathsf T}x_0,\qquad y_t=Px_t=x_t-\bar x\one.
\]
Every $W_{ij}$ is symmetric, fixes $\one$, and commutes with $P$.
The basic complete-graph identity is
\begin{equation}\label{eq:laplacian}
 \sum_{i<j}d_{ij}d_{ij}^{\mathsf T}=NP.
\end{equation}

\begin{theorem}[Mean and sharp energy law]\label{thm:energy}
The sample mean is invariant on every realization.  With
\begin{equation}\label{eq:mu-rho}
 \mu=1-\frac{2\eta}{N-1},\qquad
 \rho=1-\frac{4\eta(1-\eta)}{N-1},
\end{equation}
one has, for every $t\geq0$,
\begin{equation}\label{eq:first-energy}
 \EE y_t=\mu^t y_0,
 \qquad
 \EE\|y_t\|^2=\rho^t\|y_0\|^2.
\end{equation}
\end{theorem}

\begin{proof}
Equation~\eqref{eq:update} fixes $\one^{\mathsf T}x$ pathwise.  Since
$K_N$ has $N(N-1)/2$ edges, averaging~\eqref{eq:update} and using
\eqref{eq:laplacian} gives
\[
 \EE W=I-\frac{2\eta}{N-1}P,
\]
which acts by $\mu$ on $\one^\perp$.  Independence of the new edge from the
past proves the first identity by iteration.  Since
$(d_{ij}d_{ij}^{\mathsf T})^2=2d_{ij}d_{ij}^{\mathsf T}$,
\[
 W_{ij}^2=I-2\eta(1-\eta)d_{ij}d_{ij}^{\mathsf T}.
\]
Conditioning on $y_t$, applying~\eqref{eq:laplacian}, and iterating gives the
second identity.  It is an equality for every initial vector, hence its
rate is sharp whenever $y_0\ne0$.
\end{proof}

The phrase \emph{sharp energy owner} records the content added in the
original paper round: no spectral conclusion is needed for
Theorem~\ref{thm:energy}.

\ifnum\CRevisionRound>0
\section{Three invariant covariance blocks}
Let
\[
 \mathscr H_N=\{A\in\mathbb R^{N\times N}:A=A^{\mathsf T},\ PAP=A\}
\]
with Frobenius inner product.  For $N\geq3$ define
\begin{align*}
 \mathscr V_0&=\operatorname{span}\{P\},\\
 \mathscr V_1&=\{P\operatorname{diag}(u)P:\one^{\mathsf T}u=0\},\\
 \mathscr V_2&=\{A\in\mathscr H_N:\operatorname{diag}A=0\}.
\end{align*}
For $N=3$, $\mathscr V_2=\{0\}$.  For $N=2$, only
$\mathscr V_0=\mathscr H_2$ remains.

\begin{lemma}[Explicit orthogonal projectors]\label{lem:projectors}
For $N\geq3$ and $A\in\mathscr H_N$, set
\begin{align}\label{eq:projectors}
 \Pi_0A&=\frac{\operatorname{tr}A}{N-1}P,\qquad R=A-\Pi_0A,\nonumber\\
 u&=\frac{N}{N-2}\operatorname{diag}R,\qquad
 \Pi_1A=P\operatorname{diag}(u)P,\qquad
 \Pi_2A=A-\Pi_0A-\Pi_1A.
\end{align}
Then $\Pi_r$ are the orthogonal projectors onto $\mathscr V_r$, and
\begin{equation}\label{eq:dimensions}
 \dim\mathscr V_0=1,\qquad \dim\mathscr V_1=N-1,\qquad
 \dim\mathscr V_2=\frac{N(N-3)}2.
\end{equation}
\end{lemma}

\begin{proof}
The first residual has trace zero, so $\one^{\mathsf T}u=0$.  For such $u$,
\[
 \operatorname{diag}(P\operatorname{diag}(u)P)
 =\frac{N-2}{N}u.
\]
Hence $\Pi_1A$ has the same diagonal as $R$, and $\Pi_2A$ has zero
diagonal.  Centered matrices with zero diagonal are Frobenius orthogonal to
both $P$ and every centered diagonal pattern.  This proves orthogonality and
the projector formulas.  The map $u\mapsto P\operatorname{diag}(u)P$ is
injective on $\one^\perp$ for $N>2$.  Subtracting dimensions from
$\dim\mathscr H_N=N(N-1)/2$ proves~\eqref{eq:dimensions}.
\end{proof}

Define the second-moment transfer
\begin{equation}\label{eq:transfer}
 \cT_\eta(A)=\frac{2}{N(N-1)}\sum_{i<j}W_{ij}AW_{ij}.
\end{equation}

\begin{theorem}[Complete second-moment spectrum]\label{thm:spectrum}
The three spaces in Lemma~\ref{lem:projectors} are orthogonal invariant
blocks of $\cT_\eta$ for every $0\leq\eta\leq1$.  The restrictions to the
present nonzero blocks are scalar, with respective values
\begin{align}
 \lambda_0&=1-\frac{4\eta(1-\eta)}{N-1},\label{eq:lambda0}\\
 \lambda_1&=1-\frac{4\eta-2\eta^2}{N-1},\label{eq:lambda1}\\
 \lambda_2&=1-\frac{4\eta}{N-1}
              +\frac{4\eta^2}{N(N-1)}.\label{eq:lambda2}
\end{align}
Absent low-dimensional blocks are omitted.  If $\eta>0$, the values attached
to the present blocks are pairwise distinct, so those blocks are precisely
the corresponding full eigenspaces.  If $\eta=0$, then $\cT_0=I$: the only
eigenvalue is one, its eigenspace is all of $\mathscr H_N$, and its
multiplicity is $N(N-1)/2$.
\end{theorem}

\begin{proof}
For centered $A$, expansion of~\eqref{eq:transfer} gives
\begin{equation}\label{eq:Texpand}
 \cT_\eta(A)=(1-\frac{4\eta}{N-1})A
 +\frac{2\eta^2}{N(N-1)}\cS(A),
 \quad
 \cS(A)=\sum_{i<j}(d_{ij}^{\mathsf T}Ad_{ij})d_{ij}d_{ij}^{\mathsf T}.
\end{equation}
If $A=P$, every scalar coefficient in $\cS$ is two, so
$\cS(P)=2NP$.  If $A=P\operatorname{diag}(u)P$ with
$\one^{\mathsf T}u=0$, then
$d_{ij}^{\mathsf T}Ad_{ij}=u_i+u_j$.  The off-diagonal and diagonal entries
give $\cS(A)=NA$.  Finally, for $A\in\mathscr V_2$,
$d_{ij}^{\mathsf T}Ad_{ij}=-2a_{ij}$; its zero row sums then give
$\cS(A)=2A$.  Substitution in~\eqref{eq:Texpand} yields the three displayed
eigenvalues.
For $\eta>0$,
\[
 \lambda_0-\lambda_1=\frac{2\eta^2}{N-1}>0,
 \qquad
 \lambda_1-\lambda_2=\frac{2\eta^2(N-2)}{N(N-1)}>0,
\]
where the second comparison is used only when $\mathscr V_2$ is present.
Together with the orthogonal decomposition, this proves the full-eigenspace
statement.  At $\eta=0$ every $W_{ij}=I$, which proves the asserted merger.
\end{proof}

\begin{theorem}[Exact second moment]\label{thm:second}
Let $A_0=y_0y_0^{\mathsf T}$ and
$M_t=\EE[y_ty_t^{\mathsf T}]$.  For $N\geq3$,
\begin{equation}\label{eq:Mt}
 M_t=\lambda_0^t\Pi_0A_0+\lambda_1^t\Pi_1A_0
       +\lambda_2^t\Pi_2A_0,
\end{equation}
where the last term is absent at $N=3$.  At $N=2$,
$M_t=\lambda_0^tA_0$.
\end{theorem}

\begin{proof}
The newly selected edge is independent of the past, hence
$M_{t+1}=\cT_\eta(M_t)$.  Apply Theorem~\ref{thm:spectrum} and iterate.
Taking traces recovers Theorem~\ref{thm:energy}, because the two trace-free
blocks do not contribute.
\end{proof}
\fi

\ifnum\CRevisionRound>1
\section{Covariance and probabilistic closure}
The disagreement second moment $M_t$ is not, in general, the statistical
covariance: $\EE y_t=\mu^ty_0$ need not vanish.  The exact covariance is
therefore
\begin{equation}\label{eq:covariance}
 \operatorname{Cov}(x_t)
 =M_t-\mu^{2t}y_0y_0^{\mathsf T}.
\end{equation}
Both terms annihilate the consensus direction.

\begin{theorem}[Tail bound and almost-sure consensus]\label{thm:as}
If $N\geq2$, $0<\eta<1$, $\varepsilon>0$, and $y_0\ne0$, then
\begin{equation}\label{eq:tail}
 \Pr\{\|y_t\|\geq\varepsilon\|y_0\|\}
 \leq\frac{\rho^t}{\varepsilon^2}.
\end{equation}
For every initial vector, including $y_0=0$, one has
$x_t\to\bar x\one$ almost surely and in mean square.
\end{theorem}

\begin{proof}
Here $0\leq\rho<1$.  Markov's inequality and
Theorem~\ref{thm:energy} give~\eqref{eq:tail}; division by $\|y_0\|$ is
legitimate under the stated hypothesis.  If $y_0=0$, every update fixes the
initial consensus vector.  Otherwise, for each positive rational
$\varepsilon$, the probabilities in~\eqref{eq:tail} are summable in $t$.
Borel--Cantelli, followed by a countable intersection, gives
$\|y_t\|\to0$ almost surely.  Mean-square convergence is already the exact
energy identity.
\end{proof}

\section{Complete boundary atlas}
For $N=1$ there is no edge to sample; by definition the process is static.
For $N=2$ the disagreement line is one-dimensional and its difference is
multiplied by $1-2\eta$ at every step.  Hence $\eta=1/2$ reaches consensus
in one update.  For $N=3$, the third block has dimension zero, so inserting
a $\lambda_2$ multiplicity would be spurious.

At $\eta=0$, all $W_{ij}$ and $\cT_0$ are the identity.  The three invariant
summands therefore merge into the eigenvalue-one eigenspace
$\mathscr H_N$, of total multiplicity $N(N-1)/2$; they must not be called
three distinct eigenspaces at this endpoint.  At $\eta=1$, each $W_{ij}$ is
the transposition of coordinates $i$ and $j$, and $\rho=1$; non-consensus
data cannot converge to consensus because their disagreement norm is
constant.  This is the precise random-transposition boundary, not an
extension of the interior convergence theorem.  Constant initial vectors
are fixed for every parameter.  Values $\eta\notin[0,1]$, nonuniform edge
laws, and noncomplete graphs lie outside the frozen owner.

\section{Evidence, collisions, and Route-A boundary}
The exact certificate contains 56 spectral rows, six projector receipts,
48 exhaustive-word rows, and 2,966 scalar leaves.  It enumerates all 4,242
edge words in its declared small-time grid.  A producer-independent checker
performs 1,392 checks, while the SymPy lane closes 350 exact identities.
Two isolated producer runs are byte-identical, and 140 repaired-hash,
parser, theorem, boundary, and evaluator attacks are rejected.  The finite
grid is regression evidence; the preceding theorems are analytic for every
declared $N$ and $\eta$.

The closest registered systems are C183, which owns the permutation-valued
random-transposition chain; C203, deterministic signed-Laplacian consensus;
C312, state-dependent Hegselmann--Krause averaging; and C322, Kac's
continuous-angle collision operator.  None owns the interior relaxed random
product and its full second-moment decomposition.

The phrase \emph{probability boundary and route firewall} records the final
revision.  Under \texttt{NO\_BAD\_EULER\_OR\_ROOT\_NUMBER}, the Route-A tuple
is
\[
 (\mathrm{A0\_FAIL},\mathrm{A1\_FAIL},\mathrm{A2\_FAIL},
  \mathrm{A3\_FAIL},\mathrm{A4\_FORMAL\_HINT}).
\]
Pair maps and the moment transfer are symmetric on natural source spaces,
but this is only a formal lift hint.  Random pair words are not deterministic
isolated prime-labelled cycles, and interaction count is not a logarithmic
prime clock.  There is no orbit zeta, target Fredholm determinant,
functional equation, Weil compression, target divisor, target-zero match,
or natural same-clock unitary quantization.  Route A is rejected and Route B
remains locked.  No target arithmetic local data, Euler factors, root number,
automorphy, or Hilbert--Polya operator is claimed.

\paragraph{AI use.}
A generative language model assisted proof organization, code scaffolding,
and manuscript drafting.  The displayed derivations, independent exact
checker, symbolic lane, hostile tests, and deterministic artifacts define
the audit record.
\fi

\begingroup
\small
\begin{thebibliography}{9}
\bibitem{BoydGossip}
S. Boyd, A. Ghosh, B. Prabhakar, and D. Shah,
``Randomized gossip algorithms,''
\emph{IEEE Transactions on Information Theory} 52 (2006), 2508--2530.
DOI: \href{https://doi.org/10.1109/TIT.2006.874516}%
{10.1109/TIT.2006.874516}.
\end{thebibliography}
The source establishes the randomized-gossip owner and network setting.
The complete-graph formulas above are derived directly for the frozen model;
no literature-priority claim is made for them.
\endgroup
\end{document}
